\section{Discussion}
\label{sec:discussion}

\subsection{Feedback Quality Is the Primary Signal}
\label{sec:feedback-quality}

The primary finding of this study is not that richer reviewer profiles produce higher scores.
The average score delta from v1 to v2 profiles across the three-puzzle ablation is +0.57---a
modest increment that would be difficult to interpret as practically meaningful on its own.
What increases with profile depth is something more valuable and less obvious: the specificity,
transferability, and operational utility of the diagnostic output.

The distinction is visible in the framework column of Table~\ref{tab:ablation}. Conditions C0
and C1 produce zero named evaluative frameworks across all three puzzles. When a C0 reviewer
identifies a problem, the diagnosis is cast in the rubric's own vocabulary: ``the extraction
mechanism is insufficiently clear.'' When a C6 reviewer identifies the same problem, it is
named as a \emph{visual grammar} failure (Dana Young's lens) or a \emph{load-bearing test}
failure (Thomas Snyder's lens)---labels that are transferable across puzzles and actionable in
isolation. The difference in practical utility is not subtle. A review that scores a puzzle
24/30 and identifies that ``the Bonus B clue requires back-solving rather than forward
deduction---this is a fairness-category failure, not a difficulty problem'' is more useful
to the constructor than a review that scores 26/30 and observes that ``the extraction could
be cleaner.'' The named mechanism is the contribution; the number is secondary.

The six domain-emergent frameworks surfaced by C6 reviews---\emph{a-ha economy},
\emph{load-bearing test}, \emph{world-model consistency}, \emph{social fabric},
\emph{perceptual-shift}, and \emph{visual grammar}---did not appear in any input text
provided to the reviewers. These frameworks were identified by the research team through
two-pass manual coding of review transcripts: a first pass using automated keyword matching
for known framework names, followed by manual review for novel or paraphrased analytical
lenses. Inter-rater reliability was not assessed; this represents a limitation on the
framework-count comparisons. Whether these frameworks are generated from the profile
structure or retrieved from the model's training prior is a confound this study cannot
fully resolve (see Section~\ref{sec:limitations}). What distinguishes them from C2
frameworks is that they are applied as diagnostic tools to specific puzzle features rather
than invoked as named criteria. The quality of named frameworks is the signal, not the
count: six domain-emergent frameworks are more practically useful than ten retrieved
principle names, irrespective of what the raw count implies.

A methodological qualification applies to this claim. The profiles encode named frameworks
attributed to documented practitioners (Thomas Snyder, Dana Young, Dan Katz); these
practitioners' vocabulary is almost certainly present in the model's training data. When a
C6 review applies the \emph{load-bearing test}, the named term may reflect profile-driven
reasoning or may reflect the model's prior over Snyder's documented vocabulary. The study
measures whether the term appears in C6 reviews and not in C0/C1 reviews---not whether it
is caused by the profile rather than by the model's knowledge of the practitioner. A control
using fictional practitioner profiles with structurally equivalent lens sections would
disentangle these contributions; this is reserved for future work.

\subsection{The Calibration Inversion: Why Rubric-Only Is Dangerous}
\label{sec:calibration-inversion}

The extended 15-puzzle corpus introduces a finding that is absent from the 3-puzzle ablation
and has significant implications for practitioners: the rubric-only condition (C1) produces
an \emph{inverted} tier hierarchy. Across the five MIT/Elite-tier puzzles, C1 yields a pass
rate of approximately 40\% (2 of 5); across the five Standard Hunt puzzles, the pass rate rises to
approximately 80\% (4 of 5). The expected ordering---elite puzzles rating higher than standard puzzles
on a domain-informed rubric---is reversed by the one condition that provides rubric structure
without reviewer identity. This pattern persists across the full corpus of 15 real puzzles
spanning three quality tiers.

The mechanism is traceable. Elite-tier hunt puzzles characteristically sacrifice Clarity and
Solvability for Elegance and Reading Reward: the difficulty is intended, and the aha moment
is withheld because withholding it is the source of its value. A bare rubric that weights all
six dimensions equally treats this trade-off as a defect rather than a design choice. Standard
Hunt puzzles, by contrast, are engineered for accessibility and score uniformly across
dimensions because they optimize for consistent competence rather than concentrated brilliance
on any single axis. The rubric, applied without contextual recalibration, rewards what it can
measure uniformly and penalizes what it cannot.

Expert identity reintroduces the missing calibration. Under C6, reviewers weight the six
dimensions differently depending on the puzzle type under evaluation. Elegance and Reading
Reward are not downweighted as such---they are contextualized: a C6 reviewer understands
that a 3/5 on Clarity in an MIT-tier puzzle represents a different signal from a 3/5 on
Clarity in an accessible feeder puzzle. The profile encodes not just evaluative criteria but
the calibration layer that converts raw dimension scores into appropriate verdicts. The C6
condition is the only condition in the ablation that passes all three AoE puzzles and produces
tier-consistent verdicts across the full 15-puzzle corpus.

This finding generalizes beyond puzzle hunt evaluation. Any creative domain where quality
trades off across multiple dimensions---where ambitious work necessarily sacrifices some
properties to achieve others---is vulnerable to the calibration inversion under rubric-only
evaluation. Practitioners who adopt structured rubrics for creative work assessment without
equally structured expert identity context risk systematically disadvantaging their most
ambitious artifacts. The rubric-only condition is, empirically, the most dangerous condition
in the ablation: it appears objective, it produces structured output, and it misleads.

\subsection{The Riven Standard as a Universal Creative Quality Criterion}
\label{sec:riven-standard}

The most consequential rubric-design finding in this paper is the effect of adding the Riven
Standard as an explicit scored dimension. Named after Rand Miller's design principle
\cite{miller2016riven}---``the puzzle IS what the field does, not an obstacle overlaid on
it''---the Riven Standard operationalizes a quality distinction that explains the gulf between
the Games Magazine tier and the hunt tiers: does the puzzle embody the essential activity of
its domain, or does it merely use the domain as content?

Adding the Riven Standard as a scored seventh dimension (C11 configuration), combined with
double-weighting Elegance and Reading Reward, achieves 24.5 percentage-point tier separation
between the MIT/Elite average and the Games Magazine average, compared to 10.9pp for the
equal-weight six-dimension baseline. This improvement substantially exceeds the gain from
double-weighting Elegance and Reading Reward alone (C9 configuration), suggesting that the
two weighted dimensions do not fully capture the variance the Riven Standard explains. The
theoretical interpretation is that Elegance and Reading Reward are proxies for the Riven
Standard criterion, but adding the criterion explicitly captures residual variance that the
proxies leave unaccounted.

The Riven Standard's formulation as a scored dimension---on a five-point scale from 1
(\emph{theme is decorative; any content could substitute}) to 5 (\emph{mechanic is
inseparable from the domain and could not exist in any other context})---proves more tractable
than it might appear: reviewers in C11 conditions produce consistent Riven Standard scores
with lower inter-reviewer variance than most other dimensions. This suggests that the
underlying quality distinction, while philosophically subtle, is empirically recognizable by
domain-informed reviewers. Experiment~4 (Section~\ref{sec:rubric-validation}) validates these
results against the full 15-puzzle corpus.

The Riven Standard criterion is not, we suggest, specific to puzzle hunt design. Any creative
domain evaluation should explicitly ask whether the work embodies the essential activity of
the tradition rather than using it as decoration. Music evaluation should ask whether the
composition IS what the musical tradition does at its core, or merely sounds like it. Game
design evaluation should ask whether the mechanic IS the genre's central activity, or is
bolted on. Interactive fiction evaluation should ask whether the interactivity IS the
narrative, or is incidental to it. In each case, the Riven Standard operationalizes the
criterion that separates work that advances a tradition from work that deploys a tradition
as scenery. The puzzle hunt domain provides a tractable testbed for this principle precisely
because the tradition's core activity---using domain knowledge as a load-bearing extraction
mechanism---is well-defined and verifiable.

\subsection{Profile Design Implications}
\label{sec:profile-implications}

The ablation characterizes the internal structure of profile effectiveness with enough
precision to support specific design recommendations. The findings warrant four conclusions
about what makes a reviewer profile functional.

First, the \textbf{review lens section} is the profile's primary functional contribution.
Conditions C4 (philosophy only, no structured lens) and C5 (lens only, no philosophy) both
demonstrate that the lens is load-bearing in a way that biography and philosophy are not when
the two are separated. C5's unique detection of the duplicate label error in Puzzle II---a
verifiable construction failure that no other condition identified---suggests that the lens
checklist concentrates evaluative attention specifically on the visual-mechanical properties
of the puzzle artifact in a way that higher-level framing does not. The lens is not merely a
summary of the philosophy; it activates a different review mode.

Second, \textbf{biography and credentials} provide the calibration context that prevents
hunt-scope evaluative questions from being applied to feeder-scope puzzles and vice versa.
The calibration inversion under C1 is partly attributable to the absence of this context:
a rubric without identity cannot distinguish an ambitious puzzle that intentionally
sacrifices accessibility from a poorly designed puzzle that fails accessibility by default.

Third, the \textbf{redundancy analysis} identifies three principles that add unique signal
beyond complete profiles: \emph{Verify the Last Mile} (character-level extraction tracing
that no profile lens performs at this granularity), \emph{Blame the Player} (solver
retrospective affect, a dimension profiles do not address because they are designer-centric),
and \emph{Snyder's Computer Test} (explicit application of the test, as distinct from
reasoning in its spirit). Eight of eleven principles tested are subsumed by well-designed
profiles. The practical recommendation: supplement C6 with exactly these three principles,
adding approximately 400 tokens of context, and omit the remainder. C3 (full principles,
no profiles) achieves zero wrong verdicts on the 15-puzzle corpus against C6's two; we
recommend C6 as the default because it produces richer diagnostic output---named frameworks,
reviewer-specific concerns, voice---that C3 cannot generate, even where C3's binary accuracy
is marginally better.

Fourth, profile length itself is not the relevant variable. The v1 profiles (approximately
50 lines each) and v2 profiles (approximately 95 lines each) differ less in length than in
the presence of named evaluative frameworks and a structured review lens. Profile upgrade
should target the specificity and operational content of the lens section, not line count as
such.

\subsection{Limitations}
\label{sec:limitations}

Several limitations bound the claims warranted by this data. First, and most fundamentally,
the reviewers in this study are AI personas, not human experts. The profiles are constructed
from publicly available writings, interviews, and published work by the named practitioners,
but the evaluations they produce have not been calibrated against those practitioners' actual
judgments on the same artifacts. The claim that C6 produces ``expert-quality'' evaluation is
an inference from the internal coherence and domain specificity of the output, not a validated
equivalence to human expert assessment. A calibration study---in which the same puzzles are
evaluated both by the AI panel and by the named human experts or their recognized proxies---is
necessary before the system's verdicts can be treated as proxies for human expert review. This
calibration study is identified as a target for future work but was not feasible within the
scope of the present paper.

Second, puzzle hunt design is a domain with unusually well-defined quality criteria, recognized
expert practitioners, and a documented evaluative vocabulary. The evidence that rich profiles
improve AI panel performance in this domain does not straightforwardly generalize to domains
with less codified quality criteria or less prominent expert figures. Adapting the profile
architecture to a new domain requires identifying its equivalent of the hunt community's
recognized practitioners and constructing profiles that encode their specific evaluative
frameworks---a non-trivial effort whose costs and returns in other domains are unknown.
Different creative domains may require different Review Lens formulations, different rubric
dimensions, or different weighting schemes to achieve comparable tier separation.

Third, the ground-truth problem is not fully resolved. The paper treats the extended-corpus
tier assignments (MIT/Elite, Standard Hunt, Games Magazine) as ground truth against which
panel verdicts can be assessed, but these assignments are themselves judgments made by the
authors based on puzzle provenance and community consensus, not independent oracle verdicts.
The verdicts reported here are AI-generated assessments of AI-generated and human-generated
puzzles; no external oracle determines which assessments are correct. The paper reports verdicts
and separation metrics, not claims that any particular verdict is definitively right.

Fourth, the 15-puzzle extended corpus, while carefully stratified across three tiers and
multiple domains, is not a random sample from any well-defined population. Selection was
guided by tier representativeness, domain variety, and the availability of puzzles with known
provenance. Conclusions about pass rates by tier and condition should be interpreted as
tendencies in a purposive sample, not as population-level statistics.

Fifth, \emph{profile construction.} The 29 reviewer profiles were constructed by the
research team based on their reading of each practitioner's published work, interviews,
and documented design positions---not based on self-descriptions by the named individuals.
The profiles synthesize attributed positions rather than directly representing any
individual's views, and this limits any claim about simulation fidelity.
